\chapter{Entropy}

\section*{Introduction}

Entropy is often called ``disorder,'' but a better metaphor is ``counting.'' A shuffled deck has high entropy not because it is disordered, but because there are vastly more ways to arrange a shuffled deck than a sorted one. Entropy counts the number of ways a system can be arranged while looking the same macroscopically. A gas spread uniformly has high entropy because there are many molecular arrangements producing the same density. A gas in one corner has low entropy---essentially only one arrangement. The second law says: systems evolve toward the macrostate with the most microstates---toward maximum entropy.


The thermodynamic definition:
\[
dS \geq \frac{\delta Q}{T}
\]
with equality for reversible processes. The statistical definition:
\[
S = k_B \ln \Omega
\]
where $k_B$ is Boltzmann's constant and $\Omega$ is the number of microstates.

\section*{Fundamental Equations Used}

The derivation starts from the fundamental postulate of statistical mechanics: all accessible microstates of an isolated system in equilibrium are equally probable.

\section*{Assumptions}

\textbf{Assumptions:} The system is in or near thermal equilibrium ($T$ well-defined). All microstates are equally probable (fundamental postulate). The system is large enough that fluctuations are negligible.


\textbf{Limitations:} Thermodynamic definition requires well-defined $T$ (breaks down far from equilibrium). Statistical definition requires $\Omega$, computationally intractable for most systems. Does not apply to systems with long-range interactions without modification.


\section*{Mathematical Derivation}

\textbf{Step 1: Define the microstate and the fundamental postulate.}

Consider a macroscopic system with fixed energy $E$, volume $V$, and particle number $N$. The system occupies a phase space of dimension $6N$ (positions and momenta of all $N$ particles). A \emph{microstate} specifies the exact position and momentum of every particle. The \emph{fundamental postulate of statistical mechanics} (the equal a priori probability postulate) states: for an isolated system in equilibrium, all accessible microstates are equally probable.

\textbf{Step 2: Count the number of microstates.}

Let $\Omega(E, V, N)$ be the number of microstates accessible to the system (the number of distinct phase-space configurations consistent with the macroscopic constraints). For example, for an ideal gas of $N$ particles in volume $V$ with total energy $E$:
\begin{equation}
\Omega(E, V, N) = \frac{1}{N!}\frac{V^N}{h^{3N}}\frac{(2\pi m E)^{3N/2}}{\Gamma(3N/2 + 1)}
\end{equation}
where $h$ is Planck's constant, $m$ is the particle mass, and the $N!$ accounts for indistinguishability.

\textbf{Step 3: Define Boltzmann entropy.}

Boltzmann's entropy is defined as:
\begin{equation}
S = k_B \ln \Omega
\end{equation}
The logarithm ensures that entropy is \emph{extensive}: for two independent subsystems with $\Omega_1$ and $\Omega_2$ accessible states, the total has $\Omega_{12} = \Omega_1 \Omega_2$, and:
\begin{equation}
S_{12} = k_B \ln(\Omega_1\Omega_2) = k_B\ln\Omega_1 + k_B\ln\Omega_2 = S_1 + S_2
\end{equation}
A logarithmic function is the unique function satisfying $f(ab) = f(a) + f(b)$ for positive $a, b$.

\textbf{Step 4: Connect to the thermodynamic definition.}

For a reversible process, the heat absorbed is $dQ_{\text{rev}} = T\,dS$. Rearranging:
\begin{equation}
dS = \frac{dQ_{\text{rev}}}{T}
\end{equation}
This is the thermodynamic definition of entropy. For an irreversible process, more heat is dissipated than reversibly, so:
\begin{equation}
dS \geq \frac{\delta Q}{T}
\end{equation}
with equality for reversible processes. This is the Clausius inequality (the second law).

\textbf{Step 5: Verify the connection between the two definitions.}

For the ideal gas, compute $\Omega$ using the Sackur--Tetrode equation:
\begin{equation}
S = k_B\ln\Omega = Nk_B\left[\ln\left(\frac{V}{N}\left(\frac{4\pi m E}{3Nh^2}\right)^{3/2}\right) + \frac{5}{2}\right]
\end{equation}
Differentiating with respect to energy at constant $V$ and $N$:
\begin{equation}
\left(\frac{\partial S}{\partial E}\right)_{V,N} = \frac{3Nk_B}{2E} = \frac{1}{T}
\end{equation}
which matches the thermodynamic definition since $E = \frac{3}{2}Nk_BT$ for an ideal gas. This confirms the two definitions are equivalent.

\textbf{Step 6: State the second law.}

The second law of thermodynamics states that for an isolated system:
\begin{equation}
\boxed{S = k_B\ln\Omega \text{ always increases or remains the same: } \Delta S \geq 0}
\end{equation}
This is because the system evolves toward the macrostate with the largest $\Omega$---the most probable macrostate---overwhelmingly so for large $N$. The ``arrow of time'' is the direction of increasing $\Omega$.

\section*{Characteristics of the Equation}

\begin{itemize}
\item \textbf{Extensive:} $S(A + B) = S(A) + S(B)$ for non-interacting systems.
\item \textbf{Irreversibility:} $dS \geq \delta Q/T$ means entropy increases in irreversible processes.
\item \textbf{Statistical interpretation:} $S = k_B \ln \Omega$ connects macroscopic thermodynamics to microscopic physics.
\item \textbf{Information connection:} Shannon's $H = -\sum p_i \log_2 p_i$ has the same form as Boltzmann's entropy.
\item \textbf{Heat death:} If the universe is closed, entropy increases until all gradients vanish.
\end{itemize}
\section*{Solution Methods}

\begin{enumerate}
\item \textbf{Calorimetric:} $\Delta S = \int \delta Q_{\text{rev}} / T$ for reversible processes.
\item \textbf{Statistical mechanics:} Compute partition function $Z$: $S = k_B \ln Z + k_B T (\partial \ln Z / \partial T)_V$.
\item \textbf{Third Law:} $S(T) = \int_0^T (C_p / T') \, dT'$ from absolute zero.
\item \textbf{Information-theoretic:} $S = -k_B \sum_i p_i \ln p_i$ from the probability distribution over microstates.
\end{enumerate}
\section*{Verifications}

\begin{itemize}
\item \textbf{Irreversibility:} All observed processes increase total entropy. No verified exception.
\item \textbf{Maxwell's demon:} Thought experiments resolved---demon must erase information, increasing entropy elsewhere.
\item \textbf{Black hole thermodynamics:} Bekenstein--Hawking formula consistent with black hole mechanics.
\item \textbf{Information entropy:} Shannon entropy matches Boltzmann form, confirming the deep connection.
\end{itemize}
\section*{Applications}

\begin{itemize}
\item \textbf{Heat engines:} Entropy analysis determines maximum efficiency.
\item \textbf{Mixing:} Entropy of mixing quantifies irreversibility of combining substances.
\item \textbf{Chemical reactions:} $\Delta G = \Delta H - T\Delta S$ determines equilibrium constants.
\item \textbf{Black holes:} Bekenstein--Hawking $S = k_B A / (4 \ell_P^2)$ connects thermodynamics to quantum gravity.
\item \textbf{Information theory:} Shannon entropy bounds information transmission rates.
\end{itemize}
\section*{What Richard Feynman Would Have Asked}

\subsection*{1. Plain-English Translation}
\textit{``What is entropy, really? Not the textbook definition---what is it?''}

Entropy counts the number of ways you can rearrange microscopic details without changing what you see macroscopically. A gas filling a room uniformly has enormous entropy because there are unimaginably many ways to distribute molecules that all look the same. A gas in one corner has low entropy because there is essentially only one way. The second law says: systems evolve from less probable (low entropy) to more probable (high entropy) simply because there are more ways to be in the high-entropy state.

The ``arrow of time'' is entropy. The past had lower entropy than the present (the Big Bang was very low entropy), and the future will have higher. This is why we remember the past and not the future, why eggs break but do not unbreak, why cream mixes into coffee but does not unmix. Entropy is the physical quantity distinguishing past from future.

\subsection*{2. Localized Mechanics}
\textit{``What is happening at the molecular level when entropy increases?''}

At the molecular level, entropy increase is just molecules exploring more of the available space. When a gas expands, molecules are not ``pushed'' outward---they are simply more likely in the larger volume because there are more arrangements there. When heat flows from hot to cold, fast molecules collide with slow ones, transferring energy until all have roughly the same average energy (equilibrium). The number of microstates at equilibrium is vastly larger.

He would press you on ``irreversible'': from the molecular perspective, every collision is reversible (time-reversal symmetric). So why is macroscopic process irreversible? Because the reverse (all molecules spontaneously returning to the hot end) is not impossible---just overwhelmingly improbable. The probability is $\sim 2^{-N}$ where $N \sim 10^{23}$---effectively zero.

\subsection*{3. Testing Extreme Limits}
\textit{``What happens to entropy at extreme conditions?''}

At absolute zero, the entropy of a perfect crystal is zero (Third Law): exactly one microstate (ground state), so $S = k_B \ln 1 = 0$. At the Big Bang, the observable universe had very low entropy---this initial low entropy gives the universe its arrow of time.

In a black hole, entropy is proportional to the event horizon area, not volume: $S = k_A / (4 \ell_P^2)$. This suggests information content of a region scales with boundary area, not volume---the holographic principle, one of modern physics' most profound insights.

\subsection*{4. Dimensionality and Geometry}
\textit{``Does the shape or size of a system affect its entropy?''}

Yes. Entropy of a gas scales with volume: $S \propto N \ln V$ (Sackur--Tetrode equation). Doubling volume increases entropy by $N \ln 2$. In confined geometries (thin films, narrow channels), available phase space is reduced and entropy is lower. Shape matters because it determines the density of states.

In black hole physics, entropy is proportional to event horizon area, not enclosed volume---the opposite of normal thermodynamics. The resolution lies in the holographic principle: information in a volume is encoded on its boundary.

\subsection*{5. Cross-Disciplinary Connection}
\textit{``Where else does entropy appear outside of physics?''}

Entropy appears in information theory (Shannon entropy measures message uncertainty), ecology (entropy of species distributions measures biodiversity), economics (entropy of wealth distributions measures inequality), and machine learning (cross-entropy measures predicted vs.\ actual probability distributions). The mathematical universality---counting microscopic arrangements that produce the same macroscopic outcome---means entropy appears wherever uncertainty, randomness, or multiplicity of configurations matters.


%=============================================================
% Chapter 23: Euler Equation (Fluid)
%=============================================================
\section*{FAQ}

\textbf{Q1: Why does entropy always increase?}

A: Because there are vastly more disordered states than ordered ones. A system evolves toward higher entropy because there are more ways to be disordered. It is not a force---it is statistics. The probability of spontaneous decrease is not zero but astronomically small for macroscopic systems.

\textbf{Q2: Can entropy decrease locally?}

A: Yes---but only if it increases somewhere else by at least as much. A refrigerator decreases entropy inside but increases it in the room. A living organism maintains low internal entropy by consuming food and expelling waste. Total entropy of system plus surroundings always increases.

\textbf{Q3: What is the connection between entropy and information?}

A: Shannon's information entropy measures uncertainty in a message---how many possible messages could have been sent. Boltzmann's entropy measures uncertainty in the microstate---how many microstates are compatible. The forms are identical. Landauer's principle connects them physically: erasing one bit increases entropy by $k_B \ln 2$.
\section*{Important Bibliography}

\begin{enumerate}
\item Boltzmann, L., ``\"Uber die Beziehung zwischen dem zweiten Hauptsatze der mechanischen W\"armetheorie und der Wahrscheinlichkeitsrechnung,'' \textit{Wiener Berichte} 76, 373--435 (1877).
\item Pathria, R.K. and Beale, P.D., \textit{Statistical Mechanics}, 4th ed. (Academic Press, 2022), Chapter 1.
\item Callen, H.B., \textit{Thermodynamics and an Introduction to Thermostatistics}, 2nd ed. (Wiley, 1985), Chapter 7.
\item Kittel, C. and Kroemer, H., \textit{Thermal Physics}, 2nd ed. (W.H. Freeman, 1980), Chapter 6.
\end{enumerate}

